\documentclass[10pt]{beamer}
\usetheme{metropolis}
\usepackage{graphicx}
\usepackage{listings}
\usepackage{booktabs}
\usepackage{xcolor}
\usepackage{hyperref}
\usepackage{tikz}

% ── Listing styles ──────────────────────────────────────────────
\definecolor{codebg}{HTML}{F5F5F5}
\definecolor{codegreen}{HTML}{2E7D32}
\definecolor{codegray}{HTML}{757575}
\definecolor{codepurple}{HTML}{7B1FA2}
\definecolor{codeblue}{HTML}{1565C0}

\lstdefinestyle{rstyle}{
  language=R,
  backgroundcolor=\color{codebg},
  basicstyle=\ttfamily\scriptsize,
  keywordstyle=\color{codeblue}\bfseries,
  stringstyle=\color{codegreen},
  commentstyle=\color{codegray}\itshape,
  breaklines=true,
  frame=single,
  rulecolor=\color{codegray!30},
  showstringspaces=false,
  tabsize=2,
  morekeywords={library,ggplot,aes,geom_bar,geom_col,geom_text,geom_point,
    coord_flip,theme_tufte,theme_minimal,labs,scale_y_continuous,
    element_blank,element_line,fct_reorder,mutate}
}

\lstdefinestyle{pythonstyle}{
  language=Python,
  backgroundcolor=\color{codebg},
  basicstyle=\ttfamily\scriptsize,
  keywordstyle=\color{codeblue}\bfseries,
  stringstyle=\color{codegreen},
  commentstyle=\color{codegray}\itshape,
  breaklines=true,
  frame=single,
  rulecolor=\color{codegray!30},
  showstringspaces=false,
  tabsize=4,
  morekeywords={import,from,as,pd,plt,sns,np,fig,ax,set_visible}
}

% ── Metadata ────────────────────────────────────────────────────
\title{Tufte's Principles of Graphical Excellence}
\subtitle{Workshop 1 --- Module 2}
\author{Prof.Asoc.\ Endri Raco}
\institute{Data Visualization for Data Scientists \\ UNYT Tirana}
\date{2025/26}

\begin{document}

% ================================================================
\begin{frame}
  \titlepage
\end{frame}

% ================================================================
\begin{frame}{Module Roadmap}
  \begin{enumerate}
    \item The data-ink ratio: theory and application
    \item Chartjunk: a taxonomy of visual noise
    \item The lie factor: graphical integrity
    \item Small multiples and sparklines
    \item Integration of text, numbers, and graphics
  \end{enumerate}

  \begin{exampleblock}{Learning Outcome}
    You will be able to evaluate any chart against Tufte's core
    principles and systematically improve its data-ink ratio.
  \end{exampleblock}
\end{frame}

% ================================================================
\section{The Data-Ink Ratio}
% ================================================================

\begin{frame}{Tufte's Central Principle}
  \begin{center}
    \includegraphics[width=0.85\textwidth]{figures_m02/data_ink_formula}
  \end{center}

  \vspace{4pt}
  Two operational corollaries:
  \begin{enumerate}
    \item Above all else, \textbf{show the data}.
    \item Erase non-data ink, within reason.
  \end{enumerate}

  \vspace{4pt}
  \emph{Source: Tufte, E.~R. (1983). The Visual Display of Quantitative
  Information, pp.~93--96.}
\end{frame}

% ================================================================
\begin{frame}{Progressive Data-Ink Maximisation}
  \begin{center}
    \includegraphics[width=\textwidth]{figures_m02/data_ink_steps}
  \end{center}

  \begin{alertblock}{Pro-Tip}
    Apply the ``eraser'' iteratively: remove one element, check if
    anything is lost. If not, it was non-data ink.
  \end{alertblock}
\end{frame}

% ================================================================
\begin{frame}{The Range-Frame: Tufte's Axis Innovation}
  \begin{center}
    \includegraphics[width=0.92\textwidth]{figures_m02/range_frame}
  \end{center}

  \begin{exampleblock}{Data Insight}
    In a range-frame the axis line extends only from the data minimum to
    the data maximum. The axis itself becomes a data display showing
    range, not an arbitrary container.
  \end{exampleblock}
\end{frame}

% ================================================================
\begin{frame}[fragile]{Data-Ink Reduction in R}
  \begin{columns}[T]
    \begin{column}{0.52\textwidth}
\begin{lstlisting}[style=rstyle]
library(ggplot2)
library(ggthemes) # theme_tufte()

# Standard ggplot bar chart
p <- ggplot(mpg_counts,
       aes(x = fct_reorder(class, n),
           y = n)) +
  geom_col(fill = "#1565C0",
           width = 0.6) +
  geom_text(aes(label = n),
            hjust = -0.2,
            size = 3) +
  coord_flip() +
  # Apply Tufte theme: removes
  # gridlines, axis ticks, borders
  theme_tufte(base_size = 11) +
  labs(x = NULL, y = "Count",
       title = "Vehicle Class")
\end{lstlisting}
    \end{column}
    \begin{column}{0.45\textwidth}
      \includegraphics[width=\textwidth]{figures_m02/r_data_ink}
    \end{column}
  \end{columns}
\end{frame}

% ================================================================
\begin{frame}[fragile]{Data-Ink Reduction in Python}
  \begin{columns}[T]
    \begin{column}{0.52\textwidth}
\begin{lstlisting}[style=pythonstyle]
import matplotlib.pyplot as plt

fig, ax = plt.subplots()
ax.barh(cats, vals,
        color="#2E7D32",
        height=0.5)

# Direct labels — no legend needed
for i, v in enumerate(vals):
    ax.text(v + 0.8, i, str(v),
            va="center", fontsize=8)

# Tufte-style spine removal
ax.spines["top"].set_visible(False)
ax.spines["right"].set_visible(False)
ax.spines["left"].set_visible(False)
ax.tick_params(left=False)
ax.grid(axis="x", alpha=0.15)

ax.set_xlabel("Count")
plt.tight_layout()
\end{lstlisting}
    \end{column}
    \begin{column}{0.45\textwidth}
      \includegraphics[width=\textwidth]{figures_m02/py_data_ink}
    \end{column}
  \end{columns}
\end{frame}

% ================================================================
\section{Chartjunk}
% ================================================================

\begin{frame}{Tufte's Chartjunk Taxonomy}
  Chartjunk consists of visual elements that do not convey data:
  \begin{enumerate}
    \item \textbf{Moir\'{e} vibration} --- dense hatching patterns that
          flicker and distract the eye
    \item \textbf{The grid} --- heavy gridlines that compete with data
          marks for visual attention
    \item \textbf{The duck} --- decorative images and clip-art that
          overwhelm the actual data
  \end{enumerate}

  \vspace{6pt}
  \emph{``Chartjunk does not achieve the goals of its propagators.
  The overwhelming fact of data graphics is that they stand or fall
  on their content.''} --- Tufte (1983), p.~121
\end{frame}

% ================================================================
\begin{frame}{Moir\'{e} Vibration}
  \begin{center}
    \includegraphics[width=0.88\textwidth]{figures_m02/moire}
  \end{center}

  \begin{alertblock}{Pro-Tip}
    Never use crosshatching to distinguish categories. Distinct hues at
    moderate saturation are perceptually superior and accessible.
  \end{alertblock}
\end{frame}

% ================================================================
\begin{frame}{The Duck: Decoration vs.\ Data}
  \begin{center}
    \includegraphics[width=0.88\textwidth]{figures_m02/duck_vs_data}
  \end{center}

  \begin{exampleblock}{Data Insight}
    Tufte borrows ``the duck'' from architecture (Venturi et al., 1972):
    a building shaped like a duck to advertise a duck shop. In data
    graphics the decoration becomes the message, displacing the data.
  \end{exampleblock}
\end{frame}

% ================================================================
\begin{frame}{The 3-D Pie: Chartjunk in Action}
  \begin{columns}[T]
    \begin{column}{0.48\textwidth}
      \includegraphics[width=\textwidth]{figures_m02/chartjunk_3d_pie}
    \end{column}
    \begin{column}{0.48\textwidth}
      \includegraphics[width=\textwidth]{figures_m02/chartjunk_fixed}
    \end{column}
  \end{columns}

  \vspace{4pt}
  Problems with 3-D pies:
  \begin{itemize}
    \item Perspective distorts angular area (rear slices appear smaller)
    \item Shadow and explosion add non-data ink
    \item Humans compare lengths far better than angles
  \end{itemize}
\end{frame}

% ================================================================
\begin{frame}[fragile]{Removing Chartjunk in R}
  \begin{columns}[T]
    \begin{column}{0.52\textwidth}
\begin{lstlisting}[style=rstyle]
# BEFORE: default bar with grid
ggplot(df, aes(x = cat, y = val)) +
  geom_bar(stat = "identity",
           fill = "skyblue",
           color = "black")
  # heavy gridlines by default

# AFTER: Tufte-style cleanup
ggplot(df, aes(x = fct_reorder(cat,
               val), y = val)) +
  geom_col(fill = "#1565C0",
           width = 0.6) +
  coord_flip() +
  theme_minimal() +
  theme(
    panel.grid.major.y =
      element_blank(),
    panel.grid.minor =
      element_blank()
  )
\end{lstlisting}
    \end{column}
    \begin{column}{0.45\textwidth}
      \includegraphics[width=\textwidth]{figures_m02/r_data_ink}

      \vspace{4pt}
      {\small Key changes:
      flip to horizontal,
      remove y-grid,
      direct-label bars.}
    \end{column}
  \end{columns}
\end{frame}

% ================================================================
\begin{frame}[fragile]{Removing Chartjunk in Python}
  \begin{columns}[T]
    \begin{column}{0.52\textwidth}
\begin{lstlisting}[style=pythonstyle]
# BEFORE: default matplotlib bar
fig, ax = plt.subplots()
ax.bar(cats, vals)
# black border, grey background

# AFTER: systematic cleanup
fig, ax = plt.subplots()
ax.barh(cats, vals,
        color="#2E7D32", height=0.5)

# Remove 3 of 4 spines
for sp in ["top","right","left"]:
    ax.spines[sp].set_visible(False)

# Replace y-axis with direct labels
ax.tick_params(left=False)
for i, v in enumerate(vals):
    ax.text(v+0.8, i, str(v),
            va="center")

ax.grid(axis="x", alpha=0.15)
\end{lstlisting}
    \end{column}
    \begin{column}{0.45\textwidth}
      \includegraphics[width=\textwidth]{figures_m02/py_data_ink}

      \vspace{4pt}
      {\small Python requires explicit
      spine removal; R's \texttt{theme\_tufte()}
      handles it declaratively.}
    \end{column}
  \end{columns}
\end{frame}

% ================================================================
\section{Graphical Integrity}
% ================================================================

\begin{frame}{The Lie Factor}
  \begin{center}
    \Large
    $\displaystyle \text{Lie Factor} = \frac{\text{Size of effect in graphic}}{\text{Size of effect in data}}$
  \end{center}

  \vspace{6pt}
  \begin{itemize}
    \item Lie factor = 1.0 means the visual change matches the data change
    \item Values outside $[0.95,\; 1.05]$ signal distortion
    \item Most common sources: truncated axes, area-encoded bars,
          perspective projections
  \end{itemize}

  \begin{alertblock}{Pro-Tip}
    Calculate the lie factor explicitly: measure the percentage change in
    the visual encoding (height, area, position) and divide by the
    percentage change in the data. Report it alongside your chart.
  \end{alertblock}
\end{frame}

% ================================================================
\begin{frame}{Lie Factor: Detailed Example}
  \begin{center}
    \includegraphics[width=0.92\textwidth]{figures_m02/lie_factor_detail}
  \end{center}

  \begin{exampleblock}{Data Insight}
    When bubble \emph{area} encodes a 7.5\% increase but the
    \emph{visual} impression suggests 150\% growth, the lie factor is
    $150/7.5 = 20$. Bar charts from zero keep the lie factor near 1.0.
  \end{exampleblock}
\end{frame}

% ================================================================
\begin{frame}{Truncated Axes: The Most Common Lie}
  \begin{center}
    \includegraphics[width=0.88\textwidth]{figures_m02/truncated_axis}
  \end{center}

  \begin{itemize}
    \item A 2.8\% revenue increase appears ``dramatic'' on a truncated axis
    \item The same data on a zero-based axis reveals stability
    \item Context determines when a non-zero baseline is legitimate
          (e.g., temperature anomalies)
  \end{itemize}
\end{frame}

% ================================================================
\begin{frame}[fragile]{Detecting Lie Factors in R}
  \begin{columns}[T]
    \begin{column}{0.52\textwidth}
\begin{lstlisting}[style=rstyle]
# Compute lie factor for a chart
# where bar heights represent values

data_effect <- (new_val - old_val) /
               old_val * 100

# Measure visual effect:
# pixel height of new bar vs old bar
visual_effect <-
  (new_px - old_px) / old_px * 100

lie_factor <- visual_effect /
              data_effect

cat("Data change:", data_effect, "%")
cat("Visual change:",
    visual_effect, "%")
cat("Lie factor:", lie_factor)

# Rule of thumb:
# 0.95 <= lie_factor <= 1.05
\end{lstlisting}
    \end{column}
    \begin{column}{0.45\textwidth}
      \textbf{Worked example:}

      \begin{tabular}{lr}
        \toprule
        Metric & Value \\
        \midrule
        Old value    & \$982M \\
        New value    & \$1,010M \\
        Data $\Delta$    & 2.85\% \\
        \midrule
        Old bar (px) & 10 px \\
        New bar (px) & 50 px \\
        Visual $\Delta$  & 400\% \\
        \midrule
        \textbf{Lie Factor} & \textbf{140.4} \\
        \bottomrule
      \end{tabular}

      \vspace{6pt}
      {\small Truncating the y-axis from 975
      inflated a 2.85\% change to 400\%.}
    \end{column}
  \end{columns}
\end{frame}

% ================================================================
\begin{frame}[fragile]{Fixing Lie Factors in Python}
  \begin{columns}[T]
    \begin{column}{0.52\textwidth}
\begin{lstlisting}[style=pythonstyle]
import matplotlib.pyplot as plt

quarters = ["Q1","Q2","Q3","Q4"]
revenue = [982, 995, 1003, 1010]

fig, ax = plt.subplots()

# FIX 1: start y-axis from zero
ax.bar(quarters, revenue,
       color="#1565C0", width=0.5)
ax.set_ylim(0, 1200)

# FIX 2: add direct labels
for i, v in enumerate(revenue):
    ax.text(i, v + 15, f"${v}M",
            ha="center", fontsize=8)

# FIX 3: remove non-data ink
ax.spines["top"].set_visible(False)
ax.spines["right"].set_visible(False)

ax.set_ylabel("Revenue ($M)")
plt.tight_layout()
\end{lstlisting}
    \end{column}
    \begin{column}{0.45\textwidth}
      \textbf{Three-step fix:}

      \begin{enumerate}
        \item Set \texttt{ylim(0, max)} to anchor
              the baseline at zero
        \item Add direct value labels so readers
              do not estimate from the axis
        \item Remove non-data spines to maximise
              the data-ink ratio
      \end{enumerate}

      \vspace{4pt}
      {\small Result: lie factor drops from
      $\sim$140 to $\sim$1.0.}
    \end{column}
  \end{columns}
\end{frame}

% ================================================================
\section{Small Multiples \& Sparklines}
% ================================================================

\begin{frame}{Small Multiples}
  \begin{center}
    \includegraphics[width=0.92\textwidth]{figures_m02/small_multiples}
  \end{center}

  \begin{exampleblock}{Data Insight}
    By repeating the same design structure across panels, the viewer's
    eye is trained once and then compares freely. The design drops out;
    variation in the data emerges.
  \end{exampleblock}
\end{frame}

% ================================================================
\begin{frame}[fragile]{Small Multiples in R}
  \begin{columns}[T]
    \begin{column}{0.52\textwidth}
\begin{lstlisting}[style=rstyle]
# facet_wrap creates small multiples
ggplot(df, aes(x = month,
               y = revenue)) +
  geom_line(color = "#1565C0",
            linewidth = 0.8) +
  geom_area(fill = "#1565C0",
            alpha = 0.08) +
  facet_wrap(~region, ncol = 3,
             scales = "fixed") +
  theme_minimal(base_size = 9) +
  theme(
    strip.text =
      element_text(face = "bold"),
    panel.grid.minor =
      element_blank()
  ) +
  labs(x = NULL, y = "Revenue ($K)")
\end{lstlisting}
    \end{column}
    \begin{column}{0.45\textwidth}
      \textbf{Design rules:}

      \begin{itemize}
        \item Use \texttt{scales = "fixed"} so
              panels share the same y-axis
        \item Facet label replaces a legend
        \item Minimal grid: y-axis only
        \item Consistent colour across panels
      \end{itemize}
    \end{column}
  \end{columns}
\end{frame}

% ================================================================
\begin{frame}[fragile]{Small Multiples in Python}
  \begin{columns}[T]
    \begin{column}{0.52\textwidth}
\begin{lstlisting}[style=pythonstyle]
import matplotlib.pyplot as plt

fig, axes = plt.subplots(
    2, 3, figsize=(10, 5),
    sharey=True  # fixed y-scale
)

for ax, region in zip(
    axes.flatten(), regions):

    ax.plot(months, data[region],
            color="#1565C0", lw=1.5)
    ax.fill_between(
        months, data[region],
        alpha=0.08, color="#1565C0")
    ax.set_title(region,
                 fontsize=9,
                 fontweight="bold")

    # Tufte cleanup
    ax.spines["top"].set_visible(False)
    ax.spines["right"].set_visible(False)

plt.tight_layout()
\end{lstlisting}
    \end{column}
    \begin{column}{0.45\textwidth}
      \textbf{Python equivalent:}

      \begin{itemize}
        \item \texttt{sharey=True} ensures shared
              scale (like \texttt{scales="fixed"})
        \item Loop over axes and data
        \item Spine removal must be explicit
              per subplot
        \item \texttt{tight\_layout()} prevents
              label overlap
      \end{itemize}
    \end{column}
  \end{columns}
\end{frame}

% ================================================================
\begin{frame}{Sparklines: Word-Sized Graphics}
  \begin{center}
    \includegraphics[width=0.75\textwidth]{figures_m02/sparklines}
  \end{center}

  \begin{exampleblock}{Data Insight}
    Tufte introduced sparklines as ``intense, simple, word-sized
    graphics'' (2006, p.~7). They embed temporal context directly in
    text or tables, eliminating the need for separate chart real-estate.
  \end{exampleblock}
\end{frame}

% ================================================================
\begin{frame}[fragile]{Sparklines in R}
  \begin{columns}[T]
    \begin{column}{0.52\textwidth}
\begin{lstlisting}[style=rstyle]
# Using ggplot2 for sparkline panels
library(ggplot2)

ggplot(weekly_data,
       aes(x = week, y = value)) +
  geom_line(color = "#1565C0",
            linewidth = 0.6) +
  # Mark endpoints
  geom_point(
    data = endpoints,
    aes(color = type),
    size = 2) +
  scale_color_manual(
    values = c(min = "#E53935",
               max = "#2E7D32",
               end = "#1565C0")) +
  facet_wrap(~metric, ncol = 1,
             scales = "free_y") +
  theme_void() +  # total erasure
  theme(strip.text =
    element_text(hjust = 0))
\end{lstlisting}
    \end{column}
    \begin{column}{0.45\textwidth}
      \textbf{\texttt{theme\_void()}} erases:
      \begin{itemize}
        \item All axes and ticks
        \item All gridlines
        \item All background elements
      \end{itemize}

      \vspace{4pt}
      Only the data line and three marker
      points remain: minimum (red), maximum
      (green), and current value (blue).

      \vspace{4pt}
      {\small Alternative: \texttt{sparkline}
      package for inline HTML sparklines
      in RMarkdown/Quarto.}
    \end{column}
  \end{columns}
\end{frame}

% ================================================================
\begin{frame}[fragile]{Sparklines in Python}
  \begin{columns}[T]
    \begin{column}{0.52\textwidth}
\begin{lstlisting}[style=pythonstyle]
import matplotlib.pyplot as plt
import numpy as np

fig, axes = plt.subplots(
    4, 1, figsize=(5, 4))

for ax, metric in zip(
    axes, metrics):
    data = series[metric]

    ax.plot(data, color="#1565C0",
            linewidth=1.0)

    # Mark min, max, endpoint
    ax.plot(np.argmin(data),
            data.min(), "o",
            color="#E53935", ms=4)
    ax.plot(np.argmax(data),
            data.max(), "o",
            color="#2E7D32", ms=4)

    # Strip ALL non-data elements
    ax.set_xticks([])
    ax.set_yticks([])
    for sp in ax.spines.values():
        sp.set_visible(False)

    ax.set_ylabel(metric,
        rotation=0, labelpad=50)
\end{lstlisting}
    \end{column}
    \begin{column}{0.45\textwidth}
      \textbf{Design pattern:}

      \begin{enumerate}
        \item Plot the line (data only)
        \item Add exactly 3 marker points
        \item Remove \emph{all} spines, ticks,
              gridlines
        \item Use \texttt{set\_ylabel} with
              \texttt{rotation=0} for row labels
      \end{enumerate}

      \vspace{4pt}
      {\small Data-ink ratio approaches 1.0:
      every pixel is either the trend line
      or a key data point.}
    \end{column}
  \end{columns}
\end{frame}

% ================================================================
\section{Integration}
% ================================================================

\begin{frame}{Words and Graphics Belong Together}
  \begin{center}
    \includegraphics[width=0.82\textwidth]{figures_m02/annotated_scatter}
  \end{center}

  \begin{exampleblock}{Data Insight}
    Tufte argues that segregating text from graphics is ``a legacy of
    nineteenth-century engraving.'' Modern tools can seamlessly embed
    annotations, callouts, and contextual labels directly onto the
    data plane.
  \end{exampleblock}
\end{frame}

% ================================================================
\begin{frame}[fragile]{Annotation in R}
  \begin{columns}[T]
    \begin{column}{0.52\textwidth}
\begin{lstlisting}[style=rstyle]
ggplot(df, aes(x = spend,
               y = roi)) +
  geom_point(alpha = 0.5) +

  # Direct annotation
  annotate("text",
    x = 85, y = 65,
    label = "Top performer",
    fontface = "bold",
    color = "#2E7D32",
    size = 3) +

  # Arrow pointing to data
  annotate("segment",
    x = 83, xend = 78,
    y = 63, yend = 58,
    arrow = arrow(length =
      unit(0.2, "cm")),
    color = "#2E7D32") +

  # ggrepel for non-overlapping
  # geom_text_repel(aes(label))

  theme_minimal()
\end{lstlisting}
    \end{column}
    \begin{column}{0.45\textwidth}
      \textbf{R annotation toolkit:}

      \begin{itemize}
        \item \texttt{annotate("text")} for
              free-floating labels
        \item \texttt{annotate("segment")} for
              arrows and connectors
        \item \texttt{ggrepel} package for
              automatic label placement
        \item \texttt{ggtext} for markdown-
              formatted annotations
      \end{itemize}
    \end{column}
  \end{columns}
\end{frame}

% ================================================================
\begin{frame}[fragile]{Annotation in Python}
  \begin{columns}[T]
    \begin{column}{0.52\textwidth}
\begin{lstlisting}[style=pythonstyle]
fig, ax = plt.subplots()
ax.scatter(spend, roi, alpha=0.5)

# Annotate with arrow
ax.annotate(
    "Top performer\nROI=65%",
    xy=(78, 58),       # data point
    xytext=(85, 68),   # label pos
    fontsize=8,
    fontweight="bold",
    color="#2E7D32",
    arrowprops=dict(
        arrowstyle="->",
        color="#2E7D32",
        linewidth=1.2
    )
)

# Text box with background
ax.text(20, 55,
    "Quadrant analysis",
    fontsize=8,
    bbox=dict(
        boxstyle="round",
        facecolor="#E3F2FD",
        alpha=0.8))
\end{lstlisting}
    \end{column}
    \begin{column}{0.45\textwidth}
      \textbf{Python annotation toolkit:}

      \begin{itemize}
        \item \texttt{ax.annotate()} with
              \texttt{arrowprops} for arrows
        \item \texttt{ax.text()} with \texttt{bbox}
              for text boxes
        \item \texttt{adjustText} package for
              automatic repulsion
        \item \texttt{FancyArrowPatch} for custom
              arrow styles
      \end{itemize}
    \end{column}
  \end{columns}
\end{frame}

% ================================================================
\section{Synthesis}
% ================================================================

\begin{frame}{Tufte's Five Core Principles --- Summary}
  \begin{center}
    \includegraphics[width=0.82\textwidth]{figures_m02/tufte_summary}
  \end{center}
\end{frame}

% ================================================================
\begin{frame}{The Checklist Applied}
  For every chart you produce, ask:
  \begin{enumerate}
    \item \textbf{Data-Ink}: Can I remove any element without losing
          information?
    \item \textbf{Chartjunk}: Are there 3-D effects, moir\'{e} patterns,
          or decorative elements?
    \item \textbf{Lie Factor}: Does the visual change proportionally
          match the data change?
    \item \textbf{Comparison}: Would small multiples reveal more than
          a single cluttered panel?
  \end{enumerate}

  \begin{alertblock}{Pro-Tip}
    ``Above all else, show the data'' does not mean ``show
    \emph{only} the data.'' Annotations, labels, and titles are
    data-adjacent ink that aids interpretation.
  \end{alertblock}
\end{frame}

% ================================================================
\begin{frame}{Key Takeaways}
  \begin{enumerate}
    \item The \textbf{data-ink ratio} provides a quantitative target:
          maximise data ink, erase decoration.
    \item \textbf{Chartjunk} --- moir\'{e}, ducks, heavy grids ---
          actively harms comprehension.
    \item The \textbf{lie factor} is a measurable diagnostic: values
          outside $[0.95, 1.05]$ indicate deception.
    \item \textbf{Small multiples} and \textbf{sparklines} are Tufte's
          most practical inventions --- both available natively in
          R (\texttt{facet\_wrap}, \texttt{theme\_void}) and Python
          (\texttt{subplots}, spine removal).
  \end{enumerate}
\end{frame}

% ================================================================
\begin{frame}{Essential Readings for This Module}
  \begin{itemize}
    \item Tufte, E.~R. (1983). \emph{The Visual Display of Quantitative
          Information}, Chapters 4--6 (Data-ink, Chartjunk, Data Density).
    \item Tufte, E.~R. (2006). \emph{Beautiful Evidence}, Chapter 2
          (Sparklines).
    \item Wainer, H. (1984). ``How to Display Data Badly,''
          \emph{The American Statistician}, 38(2), 137--147.
    \item Few, S. (2004). ``Tapping the Power of Visual Perception,''
          \emph{Intelligent Enterprise}.
  \end{itemize}

  \vspace{6pt}
  \textbf{Next Module}: Visual Perception \& Pre-Attentive Attributes (W01-M03)
\end{frame}

\end{document}
